\section{A signed bound for the row determinant}\label{sec:signed}

This section proves Proposition~\ref{prop:signed-goal}, which completes the proof of Theorem~\ref{thm:real}.
In the case with doubled interaction OpenAI's integrand contains a determinant $\det\mathsf L$ of size $n$
that the energy method bounds by $6^n$. We prove a bound for it that is exponentially small in $N^2$ on
the near-extremal configurations. The argument is the chain
\begin{equation}\label{eq:chain}
 |\det\mathsf L|\le11^{k_+}\det\mathsf L_{--}
 =11^{k_+}\prod_{b<c}(1-x_bx_c)\prod_{b,p}(1-x_bx_p)^{-1}\cdot\frac{Z_s(\Pi_-\Pi_+)}{Z_s(\Pi_+)} .
\end{equation}
The inequality comes from a Pick matrix, a Schur complement and a contour bound
(\S\ref{ssec:pick}--\ref{ssec:T}). The identity is Heine's formula (\S\ref{ssec:heine}). The last factor is a
ratio of two log-gas integrals. We bound the numerator from above and the denominator from below
(\S\ref{ssec:loggas}). Positive nodes close to $1$ need a separate argument (\S\ref{ssec:near1}). An
inequality for the logarithmic energy of signed measures keeps the combined estimate concave
(\S\ref{ssec:reflection}). Prescribed columns add prescribed rows (\S\ref{ssec:fixedrows}). The hypotheses of
the reduction hold on all near-extremal configurations (\S\ref{ssec:class}). The section ends with the trial
measures (\S\ref{ssec:signedcert}).

\subsection{The determinant of the row functions}\label{ssec:pick}
In the case with doubled interaction OpenAI's integrand contains the determinant of the row functions at
points $x_1,\ldots,x_n\in(-1,1)\setminus\{0\}$, which we call the nodes \cite[Prop.~5.3]{OpenAI2026}. With
$\mathsf e=n-1-2g$ it factors as
\begin{equation}\label{eq:detphi}
 \det\varphi(x)=C_{n,g}\,V(x)\prod_a|x_a|^{-(n-1-g)}\det\mathsf L(x),
\end{equation}
where $C_{n,g}$ depends only on $n$ and $g$, and $V(x)=\prod_{a<b}(x_b-x_a)$. To describe $\mathsf L$ let
\[
 B(z)=\prod_a\frac{z-x_a}{1-x_az},\qquad \cK_{ab}=\frac1{1-x_ax_b},
\]
be the Blaschke product of the nodes and the Pick matrix of the Szeg\H{o} kernels $k_a(z)=1/(1-x_az)$. Put
$w_a=1$ at a negative node and $w_a=-5$ at a positive node. The value $-5=1-6$ comes from the factor
$\frac32$ in the row functions $A_r$. With $E=\operatorname{diag}(w_aB'(x_a)x_a^{\mathsf e})$ one has
$\det\mathsf L=\det(\cK+E)/\det\cK$. In the basis $k_a$ the operator $\mathsf L$ acts on the model space
$\cK_B=H^2\ominus BH^2$, which has dimension $n$. OpenAI's argument uses only $\|\mathsf L\|\le6$, that is
$|\det\mathsf L|\le6^n$, which is negligible at the scale $e^{N^2}$. Numerical evaluation suggests that
$\log|\det\mathsf L|$ is of order $-N^2$ near the extremal configuration (Remark~\ref{rem:L}). We now prove a
bound of this order.

Let $k_-$ and $k_+$ be the numbers of negative and positive nodes. For a weight $w\in\R$ let
$F(w)=\cK+\operatorname{diag}(E_-,wE_+)$, where $E_-$ and $E_+$ are the diagonal blocks of
$\operatorname{diag}(B'(x_a)x_a^{\mathsf e})$ at the negative and the positive nodes. So
$\det\mathsf L=\det F(-5)/\det\cK$. The case $w=1$ is a Pick matrix.

\begin{lemma}\label{lem:pick}
For every configuration $0\le F(1)\le2\cK$.
\end{lemma}

\begin{proof}
For a real vector $c$ put $\phi(z)=B(z)\sum_ac_a/(z-x_a)$. This is the function of $\cK_B$ with node values
$\phi(x_a)=c_aB'(x_a)$, and $\|\phi\|_{H^2}^2=c^{\mathsf T}\cK c$. By the residue theorem on the unit disc
\[
 c^{\mathsf T}(F(1)-\cK)c=\sum_a\frac{x_a^{\mathsf e}\phi(x_a)^2}{B'(x_a)}=\frac1{2\pi i}\oint_{|z|=1}\frac{z^{\mathsf e}\phi(z)^2}{B(z)}\,dz .
\]
On the unit circle $|z^{\mathsf e}\phi^2/B|=|\phi|^2$, so the right side lies in
$[-c^{\mathsf T}\cK c,\,c^{\mathsf T}\cK c]$.
\end{proof}

\subsection{A Schur complement bound}\label{ssec:schur}
The determinant $\det F(-5)=\det(F(1)-6\operatorname{diag}(0,E_+))$ is a perturbation of a Pick matrix in the
positive block only. The following inequality bounds it by the negative block.

\begin{theorem}\label{thm:C3}
Let $M=\begin{pmatrix}M_1&R^*\\R&M_2\end{pmatrix}$ be a Hermitian block matrix with $0\le M\le\beta$ and
$M_1>0$. Let $T$ be Hermitian of the size of $M_2$, with $-\tau\le T\le\tau$, and let $\gamma\ge0$. Then
\[
 \Bigl|\det\begin{pmatrix}M_1&R^*\\R&M_2-\gamma T\end{pmatrix}\Bigr|\le(\beta+\gamma\tau)^{\dim M_2}\det M_1 .
\]
\end{theorem}

\begin{proof}
The determinant equals $\det M_1\cdot\det(S-\gamma T)$, where $S=M_2-RM_1^{-1}R^*$ is the Schur complement.
Since $M\ge0$, we have $S\ge0$, and $S\le M_2\le\beta$. So $-\gamma\tau\le S-\gamma T\le\beta+\gamma\tau$, and
every eigenvalue of $S-\gamma T$ has modulus at most $\beta+\gamma\tau$.
\end{proof}

We apply this in the metric of $\cK$, that is, after the congruence by a Cholesky factor of $\cK$. By
Lemma~\ref{lem:pick} the hypotheses hold with $\beta=2$ and $\gamma=6$, provided that $\pm E_+\le\tau\cK$ as
quadratic forms. The result is
\begin{equation}\label{eq:C3}
 |\det\mathsf L|\le(2+6\tau)^{k_+}\,\det\mathsf L_{--},\qquad
 \det\mathsf L_{--}=\frac{\det(\cK_{--}+E_-)}{\det\cK_{--}},
\end{equation}
where $\cK_{--}$ is the block of $\cK$ at the negative nodes. The factor $(2+6\tau)^{k_+}$ is $e^{O(N)}$. All
the information of order $N^2$ is in $\det\mathsf L_{--}$, which depends on the positive nodes only through
the values $B_+(x_b)$ of their Blaschke product at the negative nodes.

\subsection{A contour bound}\label{ssec:T}
It remains to bound $\tau$. The quadratic form of $E_+$ is the sum of the residues of $z^{\mathsf e}\phi^2/B$ at
the positive nodes, which lie in the right half of the disc. So we integrate over the boundary of the right
half-disc instead of the circle.

\begin{lemma}\label{lem:T}
Let $M_*=\sup_{-1<y<1}|y|^{\mathsf e}/|B(iy)|$. For every real vector $c$
\[
 \Bigl|\sum_{x_a>0}B'(x_a)x_a^{\mathsf e}c_a^2\Bigr|\le\Bigl(1+\frac{M_*}2\Bigr)c^{\mathsf T}\cK c .
\]
If $\frac{\mathsf e}n-\frac1n\sum_a\Bigl(\frac{t^2}{t^2+x_a^2}-\frac{x_a^2t^2}{1+x_a^2t^2}\Bigr)>0$ for every
$t\in(0,1]$, then $M_*=1$, and the bound is $\frac32c^{\mathsf T}\cK c$.
\end{lemma}

\begin{proof}
With $\phi$ as in the proof of Lemma~\ref{lem:pick}, the sum is the contour integral of $z^{\mathsf e}\phi^2/B$
over the boundary of the right half-disc, divided by $2\pi i$. No node lies on this boundary. On the arc the
integrand has modulus $|\phi|^2$, so the arc contributes at most $\|\phi\|^2$. On the diameter $z=iy$ the
modulus is at most $M_*|\phi(iy)|^2$. By the Fej\'er--Riesz inequality on a diameter,
$\int_{-1}^1|\phi(iy)|^2dy\le\pi\|\phi\|^2$. So the diameter contributes at most $\frac12M_*\|\phi\|^2$.

For the second statement write $\log(|y|^{\mathsf e}/|B(iy)|)=nv(y)$. Then
$v(y)=\frac{\mathsf e}n\log|y|+\frac1{2n}\sum_a\log\frac{1+x_a^2y^2}{y^2+x_a^2}$, and $v(1)=0$ because $|B(i)|=1$.
Differentiation gives $tv'(t)=\frac{\mathsf e}n-\frac1n\sum_a(\frac{t^2}{t^2+x_a^2}-\frac{x_a^2t^2}{1+x_a^2t^2})$.
If this is positive on $(0,1]$, then $v<0$ on $(0,1)$. Since $v$ is even, $M_*=1$.
\end{proof}

The formal proof avoids the residue theorem. It subtracts the poles with divided differences, integrates the
remaining holomorphic function over the half-disc with primitives, and proves the Fej\'er--Riesz step by the
same device on the upper half-disc. With $\tau=\frac32$, \eqref{eq:C3} becomes
\begin{equation}\label{eq:eleven}
 |\det\mathsf L|\le11^{k_+}\det\mathsf L_{--}.
\end{equation}
The condition on $t$ asks that the nodes are not too concentrated near $0$ at any scale. At the extremal
measure its left side is at least $0.33$ for every $t$.

\subsection{Heine's formula for the negative block}\label{ssec:heine}
For $s\ge1$ and a positive continuous function $f$ on $[-1,1]$ put
\[
 Z_s(f)=\int_{[0,\pi]^s}\prod_{l<l'}(\cos\theta_l-\cos\theta_{l'})^2\prod_lf(\cos\theta_l)\,d\theta .
\]
By Heine's formula $Z_s(f)$ is $s!$ times the Hankel determinant of the moments of $f(u)(1-u^2)^{-1/2}du$.
Let
\[
 \Pi_-(u)=\prod_{x_b<0}(1+x_b^2+2x_bu),\qquad \Pi_+(u)=\prod_{x_p>0}(1+x_p^2+2x_pu)^{-1}.
\]

\begin{theorem}\label{thm:heine}
Let $s=k_+-g$. If $s\ge1$ and $k_-\le\mathsf e+1$, then
\[
 \det\mathsf L_{--}=\prod_{b<c}(1-x_bx_c)\prod_{b,p}(1-x_bx_p)^{-1}\cdot\frac{Z_s(\Pi_-\Pi_+)}{Z_s(\Pi_+)},
\]
where $b,c$ run over the negative nodes and $p$ over the positive nodes.
\end{theorem}

\begin{proof}[Sketch]
The full proof is formalized. We sketch it for orientation. One first writes
$\det\mathsf L_{--}=\det[x_b^j+B_+(x_b)x_b^{m'-j}]_{b,j<k_-}/V(x_-)$ with $m'=\mathsf e+k_--1$, by Lagrange
interpolation at the negative nodes. Since $1+x^2+2xu=x\bigl(x+x^{-1}+2u\bigr)$, the factor $\Pi_-$ is a
polynomial in $u$ with roots $-\frac12(x_b+x_b^{-1})$. By Christoffel's formula, $Z_s(\Pi_-\Pi_+)/Z_s(\Pi_+)$
is a determinant of the monic orthogonal polynomials of the weight $\Pi_+(u)(1-u^2)^{-1/2}$ at these roots,
divided by a Vandermonde determinant. That weight is of Bernstein--Szeg\H{o} type: $\Pi_+$ is the reciprocal
of $|P_+(e^{i\theta})|^2$ for the polynomial $P_+(z)=\prod_p(1+x_pz)$. Its orthogonal polynomials of degree at
least $k_+$ are explicit. At $u=-\frac12(x+x^{-1})$ they are Laurent expressions in $x$ built from $P_+$. Since
$P_+(z)/z^{k_+}\overline{P_+}(1/z)$ evaluated at $-x_b$ is $B_+(x_b)$ up to sign, the two determinants agree
after the products of Cauchy type are collected.
\end{proof}

The formal proof works with the angle variables throughout. It needs Heine's formula, Andr\'eief's identity
and the Bernstein--Szeg\H{o} polynomials, but no integration over orthogonal groups. Both integrals in
Theorem~\ref{thm:heine} are log-gas integrals. There are $s$ particles $u_l=\cos\theta_l$ on $[-1,1]$ with
logarithmic repulsion, in the external field $\log\Pi_-+\log\Pi_+$ or $\log\Pi_+$. Each field is a sum of
logarithmic potentials $l(x,u)=\log(1+x^2+2xu)$ of the nodes.

\subsection{Bounds for the two integrals}\label{ssec:loggas}
The numerator is bounded from above by the method of OpenAI's energy estimate, now with one population. The
key identity is $\log|u-v|=-\log2-\sum_{k\ge1}\frac2kT_k(u)T_k(v)$ for $u\ne v$ in $[-1,1]$. With the moments
$M_k=\sum_lT_k(u_l)$ the pair sum is $-\binom s2\log2-\sum_k\frac1k(M_k^2-\sum_lT_k(u_l)^2)$. For any sequence
$a$ and any $\bar s>0$ one has $-M_k^2\le-2\bar sa_kM_k+\bar s^2a_k^2$, which completes the square against $a$.

\begin{lemma}\label{lem:LGup}
Let $\bar s\ge s$ and let $a$ be a summable sequence. With $\mathsf T(a,u)=\sum_ka_kT_k(u)/k$ and
$\|a\|^2=\sum_ka_k^2/k$,
\[
 \log Z_s(f)\le-s^2\log2+2\bar s^2\|a\|^2+s\sup_{[-1,1]}\bigl(\log f-4\bar s\,\mathsf T(a,\cdot)\bigr)+o(\bar s^2).
\]
\end{lemma}

The diagonal terms $\sum_lT_k(u_l)^2/k$ diverge when summed over $k$. They are regularized by a damping factor
$(1-\varepsilon_N)^k$ as in \cite[Prop.~6.1]{OpenAI2026}, which changes the bound by $o(\bar s^2)$. The completed
squares $(M_k-\bar sa_k)^2/k$ are dropped in this lemma. They are kept in \S\ref{ssec:class}.

The denominator is bounded from below by Jensen's inequality. Let $\mathsf m$ be a probability measure on
$[0,\pi]$ of the form $\mathsf m=\bigl(1+2\sum_{k\le K}\gamma_k\cos k\theta\bigr)\frac{d\theta}\pi$, with a positive
density. In the variable $u=\cos\theta$ it is a polynomial multiple of the arcsine law.

\begin{lemma}\label{lem:LGlow}
For every bounded measurable $Q$,
\[
 \log Z_s(e^Q)\ge s(s-1)\Bigl(-\log2-\sum_k\frac{2\gamma_k^2}k\Bigr)+s\int Q(\cos\theta)\,d\mathsf m(\theta)+s\,\mathrm{Ent}(\mathsf m),
\]
where $\mathrm{Ent}(\mathsf m)$ is the entropy of the density of $\mathsf m$ with respect to $d\theta$.
\end{lemma}

\begin{proof}
Apply Jensen's inequality to the product density $\prod_l\frac{d\mathsf m}{d\theta}(\theta_l)$. The logarithmic
energy of $\mathsf m$ in the variable $u$ is $-\log2-\sum_k2\gamma_k^2/k$, by the expansion of $\log|u-v|$
above.
\end{proof}

The fields of the nodes integrate against $\mathsf m$ in closed form. For $|x|<1$ one has
$\int l(x,\cos\theta)\,d\mathsf m=-2\sum_k(-x)^k\gamma_k/k$, a polynomial in $x$. The arcsine part contributes
$0$, because the logarithmic potential of the arcsine law is constant on $[-1,1]$.

\subsection{Positive nodes near $1$}\label{ssec:near1}
A positive node $x_p=1-\epsilon$ puts into $\Pi_+$ a factor $(1+x_p^2+2x_pu)^{-1}$ of size $\epsilon^{-2}$ near
$u=-1$. Lemma~\ref{lem:LGup} applies the supremum of the field to all $s$ particles. So a single such node
would change the bound by about $N\log(1/\epsilon)$, while the true change of $\det\mathsf L_{--}$ is bounded.
The following reduction removes such nodes from the log gas.

\begin{theorem}\label{thm:near1}
Let $f$ be positive and nonincreasing on $[-1,1]$, and let $w=w_r\prod_{p\in P}k_p$ with $w_r>0$,
$k_p(u)=(1+2y_pu+y_p^2)^{-1}$ and $y_p\in(0,1)$. If $|P|=2m'$ or $|P|=2m'-1$, then
\[
 \log\frac{Z_s(fw)}{Z_s(w)}\le m'\log f(-1)+\log\frac{Z_{s-m'}(fw_r)}{Z_{s-m'}(w_r)} .
\]
\end{theorem}

\begin{proof}[Sketch]
The full proof is formalized. We sketch it for orientation. For an odd number of factors add one more
factor, at a point to be chosen. Write both integrals as Hankel determinants. Moving all $y_p$ to a common
value changes the ratio monotonically. This uses a one-point comparison for the weight
$(1+2yu+y^2)^{-1}$, whose Chebyshev moments form the geometric sequence $(-y)^k$. For a common value $y$ the
$2m'$ factors behave like $m'$ particles pinned near $u=-1$. Fischer's inequality bounds the numerator by
$f(-1)^{m'}$ times the determinant of the remaining block. As $y\to1$ the Schur complement of the pinned block
in the denominator converges to the Hankel matrix of the remaining $s-m'$ particles.
\end{proof}

We apply the theorem to the positive nodes with $1-x_p<10^{-3}$. Since
$\Pi_-(-1)=\prod_b(1-x_b)^2\le4^{k_-}$, each pair of such nodes changes the bound by a factor $e^{O(N)}$. In
OpenAI's energy each such node contributes $-\beta_hn\log10^3$, which is much more. Hence a near-extremal
configuration has only $O(N)$ such nodes, and their total effect is $e^{O(N^2)}$ with a constant smaller than
the margin of the final inequality (\S\ref{ssec:signedcert}).

\subsection{Concavity: an inequality for signed logarithmic energy}\label{ssec:reflection}
We now combine OpenAI's energy with \eqref{eq:eleven}, Theorem~\ref{thm:heine} and Lemmas~\ref{lem:LGup}
and~\ref{lem:LGlow}. The total is a functional of the node measure $\sigma$, split into its negative part
$\sigma_-$ and its positive part $\sigma_+$, and of the column measure. Its quadratic part is not obviously
negative. Theorem~\ref{thm:heine} contributes the term $-\iint\log(1-xy)\,d\sigma_-(x)d\sigma_+(y)$, and
Lemma~\ref{lem:LGup} applied with a trial sequence that depends on $\sigma$ contributes a convex term. Write
\[
 I(\zeta)=-\iint\log|x-y|\,d\zeta(x)d\zeta(y),\qquad L(\zeta)=-\iint\log(1-xy)\,d\zeta(x)d\zeta(y)
\]
for real measures on $[-1,1]$ and on $[0,1]$ respectively. The second is a sum of squares of moments,
$L(\zeta)=\sum_k\frac1k\bigl(\int x^kd\zeta\bigr)^2$. The quadratic part of the combined functional is
$-\frac12I-\frac12\sum_k(B_k-2C_k)^2/k+\frac12L(\sigma_+)$, where $B_k$ and $C_k$ are the power moments of the
nodes and the columns. The term $\frac12L(\sigma_+)$ is absorbed by the following theorem.

\begin{theorem}\label{thm:reflection}
Let $\zeta=\zeta_-+\zeta_+$ be a real measure of finite logarithmic energy and total mass $0$, with $\zeta_-$ on
$[-1,0]$ and $\zeta_+$ on $[0,1]$. Let $\mathsf R$ be the reflection $x\mapsto-x$. Then
\[
 I(\zeta)\ge L(\zeta_+)+L(\mathsf R\zeta_-)\ge L(\zeta_+).
\]
\end{theorem}

\begin{proof}
Put $\zeta_1=\zeta_+$ and $\zeta_2=\mathsf R\zeta_-$, both on $[0,1]$, and let $K_1(a,b)=-\log|a-b|$ and
$K_2(a,b)=-\log(a+b)$. Since $|x-y|=a+b$ for points of opposite signs,
$I(\zeta)=K_1(\zeta_1)+K_1(\zeta_2)+2K_2(\zeta_1,\zeta_2)$, which equals
$\frac12(K_1+K_2)(\zeta_1+\zeta_2)+\frac12(K_1-K_2)(\zeta_1-\zeta_2)$.

First, $\zeta_1+\zeta_2$ has mass $0$. One has $K_1=G_{\mathbb D}+L$, where
$G_{\mathbb D}(a,b)=\log|(1-ab)/(a-b)|$ is the Green function of the unit disc, whose energy is nonnegative.
By Frullani's integral $-\log(a+b)=\int_0^\infty(e^{-(a+b)t}-e^{-t})\frac{dt}t$, so
$K_2(\zeta')=\int_0^\infty(\int e^{-at}d\zeta'(a))^2\frac{dt}t\ge0$ for $\zeta'$ of mass $0$. Hence
$(K_1+K_2)(\zeta_1+\zeta_2)\ge L(\zeta_1+\zeta_2)$.

Second, $K_1-K_2=\log\frac{a+b}{|a-b|}$ satisfies $(K_1-K_2)(\zeta')\ge L(\zeta')$ for every real $\zeta'$ on
$[0,1]$. Indeed, with $a=e^{-x}$ and $b=e^{-y}$ one has $\log\frac{a+b}{|a-b|}=\log\coth\frac{|x-y|}2
=\sum_{j\text{ odd}}\frac2je^{-j|x-y|}$. For $x,y\ge0$,
$e^{-j|x-y|}=e^{-j(x+y)}+2j\int_0^{\min(x,y)}e^{-j(x+y-2v)}dv$. Put $\Phi_j(v)=\int_{x\ge v}e^{-j(x-v)}d\zeta'$.
Then $(K_1-K_2)(\zeta')=\sum_{j\text{ odd}}\bigl(\frac2j\Phi_j(0)^2+4\int_0^\infty\Phi_j(v)^2dv\bigr)$. In
$L(\zeta')=\sum_k\frac1k(\int e^{-kx}d\zeta')^2$ the terms with odd $k$ are $\Phi_k(0)^2/k$. For even $k=j+1$,
integration by parts gives $\int e^{-(j+1)x}d\zeta'=\Phi_j(0)-\int_0^\infty e^{-(j+1)v}\Phi_j(v)\,dv$. The
Cauchy--Schwarz inequality with weights $1+\frac1j$ and $1+j$ gives
$(\int e^{-(j+1)x}d\zeta')^2\le(1+\frac1j)\Phi_j(0)^2+\frac12\int\Phi_j^2$. Summing, $L(\zeta')\le(K_1-K_2)(\zeta')$.

Together, $I(\zeta)\ge\frac12L(\zeta_1+\zeta_2)+\frac12L(\zeta_1-\zeta_2)=L(\zeta_1)+L(\zeta_2)$.
\end{proof}

By Theorem~\ref{thm:reflection} the quadratic part is at most $-\frac14I-\frac12\sum_k(B_k-2C_k)^2/k\le0$ on
differences of admissible measures. So the combined functional is concave. Its maximum is bounded by the
value at a trial measure plus the supremum of a linear functional, that is, by a constant plus the suprema
of fields of one variable. The formal proof uses a discrete form of Theorem~\ref{thm:reflection} for weighted
point configurations against a reference measure with bounded density. The prescribed rows of the next
subsection change the masses of the two kinds of nodes. To compensate, we add further points distributed by
the arcsine law, whose potential is constant on $[-1,1]$.

\subsection{Prescribed rows}\label{ssec:fixedrows}
With $m$ prescribed columns the row determinant is evaluated at the free points and at the prescribed
points. The prescribed points enter $\det\mathsf L$ as nodes but are not integrated over. Lemma~\ref{lem:T}
must then be modified. Its constant $\frac32$ need not hold when the sum runs only over the free positive
nodes while the Blaschke product and the Gram matrix contain all nodes. Already two positive nodes can
require a larger constant. We therefore keep the residues at the prescribed positive nodes as a separate
term
\[
 \mathsf W=\sum_{a\ \text{prescribed},\,x_a>0}\frac{x_a^{\mathsf e}}{(1-x_a^2)|B'(x_a)|},
\]
and bound each of its terms by the reproducing kernel estimate
$(c_aB'(x_a))^2(1-x_a^2)\le c^{\mathsf T}\cK c$.

Instead of absorbing $\mathsf W$ into $\tau$, we expand. Write $\det(F(1)+\operatorname{diag}(\cdot))$ as a sum of
principal minors over subsets $\mathcal S$ of the prescribed positive nodes. This gives
\begin{equation}\label{eq:expandS}
 |\det\mathsf L|\le11^{k_+}\sum_{\mathcal S}\Bigl(\frac6{11}\Bigr)^{|\mathcal S|}\det\mathsf L_{--}^{(\mathcal S)}\,\mathsf W_{\mathcal S},
\end{equation}
where in the term for $\mathcal S$ the nodes of $\mathcal S$ interact with the other nodes through $\log(1-xy)$
instead of $\log|x-y|$. The energy then has three prescribed populations: prescribed negative nodes,
prescribed positive nodes, and the nodes of $\mathcal S$ in \eqref{eq:expandS}. Their fields are bounded on
their domains like the field $F_1$ of Section~\ref{sec:energy}. The prescribed nodes are images of quadrature
nodes, so the mass of their measure $\sigma'$ near $1$ is limited by the spacing of the quadrature nodes:
$\sigma'([x,1])\le C_1(1-x)$ for $x\ge x_c=0.999$. For such a measure of mass $\alpha$ and a field $V$ with
$V\le M_1$ on $[0,x_c]$ and $V(x)\le c_0+r\log\frac1{1-x^2}$ on $[x_c,1)$,
\[
 \int V\,d\sigma'\le\alpha M_1+C_1\int_0^{1-x_c}\bigl(c_0+r\log(1/u)-M_1\bigr)_+\,du .
\]
The integral has a closed form. This bounds the prescribed rows without a finite optimization over their
positions.

\subsection{Near-extremal configurations}\label{ssec:class}
The reduction above needs three hypotheses on the configuration. These are the condition of
Lemma~\ref{lem:T}, a bound $O(N)$ on the number of positive nodes within $10^{-3}$ of $1$, and a fixed range for
the proportion of negative nodes. We do not need them everywhere. If OpenAI's bound at the rate $U_0$ holds
for a configuration, there is nothing to prove. Otherwise the energy of the configuration is within
$\Gamma_0N^2$ of its maximum, and we show that such configurations satisfy the three hypotheses.

For such a configuration, keep the completed squares and the nonnegative differences between the fields and
their suprema in the proof of OpenAI's estimate. Their sum is at most the deficit, a constant times
$\Gamma_0N^2$. This bounds the fraction of nodes on which a field is far below its supremum, by Chebyshev's
inequality. It also bounds every linear statistic $\sum_kF_k(M_k-\bar sa_k)/k$ of the moments, by the
Cauchy--Schwarz inequality. The condition of Lemma~\ref{lem:T} is a family of such linear statistics in $t$. It
is checked on a grid of $24$ values of $t$, with a polynomial lower bound between the grid points. The count
of nodes near $1$ follows from the field of the positive nodes near $x=1$, which lies below its supremum by a
fixed amount. The range for the proportion of negative nodes follows from polynomial minorants of the
indicator functions of the two signs.

\subsection{The trial measures}\label{ssec:signedcert}
The concave majorant of \S\ref{ssec:reflection} is bounded at a trial measure. This consists of a reference
measure for the free nodes, with a piecewise polynomial density, trial sequences for the rows, the columns
and the numerator of Theorem~\ref{thm:heine}, and the measure $\mathsf m$ of Lemma~\ref{lem:LGlow}. The bound
reads
\begin{multline*}
 \text{constant}+\alpha\sup\text{(fields of the prescribed rows)}\\
 +\sup\text{(fields of the free nodes)}\\
 +\sup\text{(column field)}\le-U_0+\tfrac12\log2+c_2\alpha .
\end{multline*}
The constant is the value at the trial measure. It consists of logarithmic energies and moments of the
reference measure, the norms of the trial sequences and the entropy of $\mathsf m$, all of which have closed
forms.

The fraction $\alpha$ of prescribed columns ranges over $[0,\frac1{16}]$. We use trial measures at $\alpha=0$,
$\frac1{200}$ and $\frac1{16}$, and interpolate linearly between consecutive ones. Every field is affine in the
interpolation parameter, so its supremum on an interval is at most the chord through its suprema at the
ends. The constant term is a quadratic polynomial in the parameter, whose curvature is an explicit
combination of the cross energies of the two end measures. It remains to bound, at each of the three trial
measures, the suprema of six fields of one variable. Appendix~\ref{app:trials} lists them and describes the
verification. The smallest margin of the resulting inequality is $0.0143$ per unit of $N^2$.

This proves Proposition~\ref{prop:signed-goal}, and with \S\ref{ssec:cases} it proves Theorem~\ref{thm:real}.
